\section{Residual certificates for remaining movement}\label{sec:residual}

A protected box bounds all remaining movement, but a good bound needs an
inner box close to the limit, and such a box is hard to find when the
iterates converge only asymptotically. This section bounds remaining
movement with quantities available at the current box instead: the
residual of one exact round, and a sensitivity bound valid on a region
that contains every later box. With a matrix sensitivity bound, the
result is an a posteriori estimate of the kind given by vector-valued
contraction principles~\cite{jachymski2016perov}. The OBBT-specific
points are the following. The order structure of OBBT removes the need
for a spectral-radius condition. The region on which the sensitivity
bound must hold can be obtained from a protected box. Only some kinds of
observed movement are valid residual information. Finally, the history of
an iteration cannot replace a verified sensitivity bound.

\subsection{Endpoint maps and invariant regions}\label{sec:residual-maps}

Represent a box $B$ by its endpoint vector $p(B)=(-\ell,u)\in\R^{2n}$, as
in Section~\ref{sec:foundations}, and write $B(p)$ for the box with
endpoint vector $p$. Then $p\le q$ componentwise exactly when
$B(p)\subseteq B(q)$, and shrinking a box decreases $p$. Fix a cutoff $U$
and let $F(p)=p(T_U(B(p)))$, defined when $K_U(B(p))\ne\emptyset$. Then
$F(p)\le p$, and monotonicity makes $F$ order preserving: $p\le q$
implies $F(p)\le F(q)$. Fixed points of $F$ are the endpoint vectors of
fixed boxes. Jacobi iteration is $p_{k+1}=F(p_k)$, a nonincreasing
sequence. The \emph{residual} at $p$ is $p-F(p)\ge0$, the movement of
the next exact round.

A set $\mathcal D$ of endpoint vectors is \emph{invariant} if $F$ is
defined on $\mathcal D$ and $F(\mathcal D)\subseteq\mathcal D$. A
sensitivity bound is useful only on an invariant region containing the
starting point, because every later iterate must stay where the bound has
been verified. Protected boxes supply such regions. For vectors
$p'\le p$, write $[p',p]=\{q:p'\le q\le p\}$.

\begin{lemma}[Invariant order intervals]\label{lem:invariant-interval}
Let $P$ be protected with threshold at most $U$, with endpoint vector
$p_P\le p_0$. Then $[p_P,p_0]$ is invariant.
\end{lemma}
\begin{proof}
For $q\in[p_P,p_0]$, the box $B(q)$ contains $P$, so
$K_U(B(q))\supseteq K_U(P)\ne\emptyset$ and $F(q)$ is defined. Order
preservation and $F(p_P)=p_P$, which is Theorem~\ref{thm:protected}(a),
give $p_P\le F(q)\le q\le p_0$.
\end{proof}

If an incumbent $\hat x$ with $f(\hat x)\le U$ lies in the current box
$B(p_0)$, the protected singleton $P=\{\hat x\}$
(Corollary~\ref{cor:original-points}) gives the region of boxes inside
$B(p_0)$ that contain $\hat x$. At a node of a search tree the incumbent
can lie outside the node box; it then supplies no such region, and
another invariant region must be found.

\subsection{A matrix majorant for the remaining movement}\label{sec:residual-tail}

\begin{theorem}[Remaining movement]\label{thm:residual}
Let $\mathcal D$ be invariant and $p_0\in\mathcal D$. Suppose a fixed
matrix $M\ge0$ satisfies
\begin{equation}\label{eq:matrix-lipschitz}
 F(p)-F(q)\le M(p-q)\qquad(p,q\in\mathcal D,\ q\le p).
\end{equation}
Let $d=p_0-F(p_0)$, let $\bar d\ge d$, and suppose a vector $e\ge0$
satisfies
\begin{equation}\label{eq:supersolution}
 \bar d+Me\le e.
\end{equation}
Then $p_k$ converges to a finite limit $p_\infty$, and for every
$m\ge0$,
\begin{equation}\label{eq:tail-bound}
 0\le p_m-p_\infty\le M^me.
\end{equation}
In particular, $p_0-p_\infty\le e$ and
$p_1-p_\infty\le Me\le e-\bar d$.
\end{theorem}
\begin{proof}
Let $\Delta_k=p_k-p_{k+1}\ge0$. Both $p_k$ and $p_{k+1}$ lie in
$\mathcal D$, and $p_{k+1}\le p_k$, so~\eqref{eq:matrix-lipschitz} gives
$\Delta_{k+1}=F(p_k)-F(p_{k+1})\le M\Delta_k$. Since $M\ge0$ preserves
componentwise inequalities, induction gives
$\Delta_k\le M^kd\le M^k\bar d$. Next we show by induction on $N\ge0$ that
\[
 \sum_{j=0}^{N-1}M^j\bar d+M^Ne\le e.
\]
For $N=0$ this reads $e\le e$. If it holds for $N$, multiplying by $M$
and adding $\bar d$ gives
$\sum_{j=0}^{N}M^j\bar d+M^{N+1}e\le\bar d+Me\le e$
by~\eqref{eq:supersolution}. Since $M^Ne\ge0$, every partial sum
$\sum_{j<N}M^j\bar d$ is at most $e$. For $N>m$,
\[
 0\le p_m-p_N=\sum_{k=m}^{N-1}\Delta_k
 \le M^m\sum_{j=0}^{N-m-1}M^j\bar d\le M^me.
\]
The sequence $(p_k)$ is nonincreasing and bounded below by $p_0-e$, so it
converges, and letting $N\to\infty$ proves~\eqref{eq:tail-bound}. For
$m=1$, inequality~\eqref{eq:supersolution} gives $Me\le e-\bar d$.
\end{proof}

Hypothesis~\eqref{eq:matrix-lipschitz} compares only ordered pairs; a
two-sided bound $|F(p)-F(q)|\le M|p-q|$ implies it. Convexity of
$\mathcal D$ is not needed, although it is useful when
\eqref{eq:matrix-lipschitz} is derived by integrating derivative bounds
along segments, as in Section~\ref{sec:constraints}. If the limit lies in
$\mathcal D$, for example when $\mathcal D$ is closed, it is a fixed
point: $0\le F(p_k)-F(p_\infty)\le M(p_k-p_\infty)\to0$, so
$F(p_\infty)=\lim_kp_{k+1}=p_\infty$.

The after-round bound $Me$ holds for an upper residual $\bar d$ as well
as for the exact one. It concerns the exact image $p_1=F(p_0)$, not an
arbitrary partially updated state. For $F(s)=s/2$, $p_0=1$, and
$\bar d=d=1/2$, the vector $e=1$ satisfies~\eqref{eq:supersolution} and
the exact tail after one round is $1/2=Me$. If no update has been
applied, however, the remaining movement is $1>e-\bar d$.

The theorem also gives a test that can refute a claimed matrix: the
exact residuals must satisfy $\Delta_k\le M^k\bar d$. An exact residual
violating this inequality shows that~\eqref{eq:matrix-lipschitz} is
false. Agreement with observed residuals does not prove it.

The next proposition identifies the best majorant obtainable from given
$M$ and $\bar d$, and gives a cheap explicit one.

\begin{proposition}[Least majorant]\label{prop:least-majorant}
Let $M\ge0$ and $r\ge0$. A finite $e\ge0$ with $r+Me\le e$ exists if
and only if the series $s=\sum_{k\ge0}M^kr$ converges. In that case $s$
satisfies $r+Ms=s$ and $s\le e$ for every such $e$. If $\rho(M)<1$, then
$s=(I-M)^{-1}r$. If a vector $w>0$ and a number $q\in[0,1)$ satisfy
$Mw\le qw$, then $\rho(M)\le q$, and
\[
 e=\frac{t}{1-q}\,w,\qquad t=\max_i\frac{r_i}{w_i},
\]
satisfies $r+Me\le e$.
\end{proposition}
\begin{proof}
If $e$ exists, the proof of Theorem~\ref{thm:residual} bounds every
partial sum of the nonnegative series by $e$, so the series converges and
$s\le e$. Conversely, passing to the limit in
$\sum_{k\le N+1}M^kr=r+M\sum_{k\le N}M^kr$ gives $s=r+Ms$. If
$\rho(M)<1$, the Neumann series gives $s=(I-M)^{-1}r$. If $Mw\le qw$ with
$w>0$, the norm $\norm{y}_w=\max_i|y_i|/w_i$ satisfies
$\norm{My}_w\le q\norm{y}_w$ for all $y$, so $\rho(M)\le q$. Finally,
$r\le tw$ gives $r+Me\le tw+\frac{tq}{1-q}w=\frac{t}{1-q}w=e$.
\end{proof}

Thus the strongest bound available from $M$ and $\bar d$ is
$s=\sum_kM^k\bar d$. When $M$ and $\bar d$ are rational and
$\rho(M)<1$, it is computed exactly by a rational linear solve. Checking
a proposed $e$ requires only the inequality~\eqref{eq:supersolution},
with no eigenvalue computation. When a protected box $P$ is also known,
the componentwise minimum of $e$ and $p_0-p_P$ bounds the remaining
movement.

\begin{example}[A noncontracting direction]\label{ex:inactive}
For the feasible set $\{0\}$ in $B_0=[-1,1]$ with $f\equiv0$ and $U=0$,
let $\phi_{[\ell,u]}(x)=0$ for $x\le u/2$ and $+\infty$ otherwise, on
every subbox $[\ell,u]$ of $B_0$. This family is valid, monotone because
$u/2$ is nondecreasing in $u$, and lower semicontinuous. For
$\ell\le0\le u$ it gives $K_0([\ell,u])=[\ell,u/2]$, so in endpoint
coordinates $F(p_1,p_2)=(p_1,p_2/2)$ on the invariant region $[0,1]^2$.
Hence~\eqref{eq:matrix-lipschitz} holds with
$M=\operatorname{diag}(1,1/2)$, whose spectral radius is $1$. From
$p_0=(1,1)$, the residual is $d=(0,1/2)$, and $e=(0,1)$
satisfies~\eqref{eq:supersolution}. The bound is exact: $p_\infty=(1,0)$.
The matrix inequality succeeds because the residual has no component in
the noncontracting direction and $M$ propagates none into it.
\end{example}

\subsection{Valid residual information}\label{sec:residual-valid}

Theorem~\ref{thm:residual} needs an upper bound $\bar d$ on the
\emph{exact} Jacobi residual. Three sources are valid: the residual of an
exact Jacobi round, for example one whose support values are proved by
primal--dual certificates (Section~\ref{sec:algorithms}); the movement of
an exact complete sequential round, which ends in a box contained in
$T_U(B)$ by Lemma~\ref{lem:sequential}(b); and the current-round ceilings
of a witness pool, which require no support solve at all. Three kinds of
observed movement are \emph{not} valid. A selective round reports no
movement in unprocessed directions; an interrupted round stops early; and
numerically safe OBBT rounds its bounds outward. Each can report less
than the exact residual.

\begin{corollary}[Residual bounds from witnesses]\label{cor:witness-residual}
Let $W\subseteq K_U(B(p_0))$ be nonempty and finite, and let $\bar d(W)$
collect its current-round ceilings: $\min_{x\in W}x_i-\ell_i$ in the
lower components and $u_i-\max_{x\in W}x_i$ in the upper components.
Then $d\le\bar d(W)$, and Theorem~\ref{thm:residual} applies with
$\bar d=\bar d(W)$.
\end{corollary}
\begin{proof}
This is Proposition~\ref{prop:current-round} written in endpoint
coordinates.
\end{proof}

A solver with a witness pool and a verified comparison matrix on an
invariant region can therefore bound all future movement of exact Jacobi
OBBT without solving a support problem. The pool may come from mixing or
from a cached frontier at a new cutoff
(Propositions~\ref{prop:mixing} and~\ref{prop:frontier}). The bounds also
apply to complete sequential rounds $S_k$ started at $B(p_0)$. By
Lemma~\ref{lem:sequential}, $S_k\subseteq B(p_k)$ and the $S_k$ have the
same limit $B(p_\infty)$, so after $m$ sequential rounds the remaining
movement is at most $p(S_m)-p_\infty\le p_m-p_\infty\le M^me$.

\subsection{A rebuilt relaxation with a verified comparison constant}\label{sec:residual-example}

The following example has McCormick coefficients that change in every
round. The theorem needs no fixed constraint matrix, only a comparison
constant valid over the invariant region.

\begin{example}[Heron's iteration from a square relaxation]\label{ex:heron}
Minimize $x^2$ over $[0,u_0]$ with the square relaxation of the
reference family and cutoff $U=r^2$, where $0<r\le u_0$. On $[0,u]$ with
$u\ge r$, the rows are $s\ge0$, $s\ge2ux-u^2$ and $s\le ux$, together
with $s\le r^2$. The upper tangent and the cutoff give
$x\le G(u)=(u^2+r^2)/(2u)$. The point $(G(u),r^2)$ attains this bound:
$r\le G(u)\le u$ because $u^2+r^2\ge2ur$ and $r\le u$; the tangent holds
with equality; and $r^2\le uG(u)=(u^2+r^2)/2$. The point $(0,0)$ keeps the
lower endpoint at $0$. The upper endpoint therefore follows Heron's
iteration $u_{k+1}=G(u_k)$, which is Newton's method for $u^2=r^2$. Since
\[
 G(u)-r=\frac{(u-r)^2}{2u}\qquad(u\ge r),
\]
the errors $u_k-r$ are nonnegative and nonincreasing, they tend to $0$,
and $(u_{k+1}-r)/(u_k-r)^2\to1/(2r)$ while $u_k>r$: the convergence is
quadratic.

The box $[0,r]$ is protected at cutoff $r^2$ by the lifted witnesses
$(0,0)$ and $(r,r^2)$; on $[0,r]$ the tangent at $r$ and the secant force
$s=r^2$ at $x=r$. Lemma~\ref{lem:invariant-interval} gives the invariant
region $\{(0,u):r\le u\le u_0\}$. There
$G'(u)=(1-r^2/u^2)/2\in[0,q]$ with $q=(1-r^2/u_0^2)/2<1/2$, so the mean
value theorem gives~\eqref{eq:matrix-lipschitz} with
$M=\operatorname{diag}(0,q)$. For $r=1$ and $u_0=2$, the residual is
$d=(0,3/4)$ and $q=3/8$, so $e=(I-M)^{-1}d=(0,6/5)$. The actual total
movement is $1$. After the first round, $Me=(0,9/20)$, while the actual
remaining movement is $1/4$. The exact iterates are $5/4$, $41/40$,
$3281/3280$; the linear bound is conservative because the convergence is
quadratic.
\end{example}

In this example the protected box $[0,r]$ equals the limit and gives the
exact remaining movement $u_0-r$ directly. The residual certificate is
useful in the opposite situation: when no inner box close to the limit is
known, but a comparison constant can be verified on a region between a
known inner box, such as an incumbent singleton, and the current box.

\subsection{Cutoff decreases}\label{sec:residual-cutoff}

When the incumbent improves, an old residual certificate does not
automatically transfer to the new cutoff, because both the map and its
residual change. A new certificate needs a new residual bound, which
witnesses supply without support solves.

\begin{corollary}[Remaining movement after a cutoff decrease]\label{cor:cutoff-tail}
Let $p$ be the endpoint vector of the current box, let the cutoff
decrease to $V$, and let $F_V$ be the endpoint map at cutoff $V$. Let
$\mathcal D\ni p$ be invariant for $F_V$, and let $M$
satisfy~\eqref{eq:matrix-lipschitz} for $F_V$ on $\mathcal D$. If $\bar d$
bounds $p-F_V(p)$ from above, for example as the ceilings at cutoff $V$
of a pool obtained by mixing or from a cached frontier, and $\bar d+Me\le
e$, then exact Jacobi OBBT at cutoff $V$ moves the endpoints of the
current box by at most $e$ in total.
\end{corollary}
\begin{proof}
Apply Theorem~\ref{thm:residual} to $F_V$.
\end{proof}

If the new incumbent lies in the current box, it supplies an invariant
region through Lemma~\ref{lem:invariant-interval}. The corollary follows
one specific iteration and needs no spectral condition. A comparison
between fixed points at two cutoffs is a different statement, and the
following proof of it uses a spectral condition.

\begin{proposition}[Shift of a fixed point]\label{prop:fixed-point-shift}
Let $V\le U$. Let $p_U$ be a fixed point of $F_U$ and $p_V\le p_U$ a
fixed point of $F_V$, both in a region $\mathcal D$ on which $F_V$
satisfies~\eqref{eq:matrix-lipschitz} with $\rho(M)<1$. If $h\ge0$
satisfies $F_U(p_U)-F_V(p_U)\le h(U-V)$, then
\[
 0\le p_U-p_V\le(I-M)^{-1}h\,(U-V).
\]
\end{proposition}
\begin{proof}
Let $x=p_U-p_V\ge0$ and $c=h(U-V)$. Then
$x=\bigl(F_U(p_U)-F_V(p_U)\bigr)+\bigl(F_V(p_U)-F_V(p_V)\bigr)\le c+Mx$.
Iterating gives $x\le\sum_{j<N}M^jc+M^Nx$ for every $N$, and $M^N\to0$.
\end{proof}

The sensitivity bound $h$ is needed only at the old fixed point, not
uniformly over a region. The condition $\rho(M)<1$ cannot simply be
dropped from this proposition. For $F_U=F_V=F$ with $F(s,t)=(s,t/2)$,
$M=\operatorname{diag}(1,1/2)$ and $h=0$, both $(1,0)$ and $(0,0)$ are
fixed points, ordered as required, yet they differ.

\subsection{What observed progress cannot certify}\label{sec:residual-history}

An exact observed residual is valid input for Theorem~\ref{thm:residual}.
The comparison matrix is different: neither it nor any nontrivial bound
on the remaining movement can be inferred from a finite strictly
decreasing trajectory alone.

\begin{example}[An observed ratio is not a comparison constant]\label{ex:ratio}
Let $F(s)=3s/4$ for $0\le s\le1/2$ and $F(s)=s/4+1/4$ for $1/2\le s\le1$.
The map is continuous and nondecreasing with $0\le F(s)\le s$, and it
is the upper-endpoint map of a valid monotone family constructed as in
Proposition~\ref{prop:history} below. From $s_0=1$ the iterates are
$s_1=1/2$ and $s_2=3/8$, with observed decrement ratio $1/4$. A geometric
extrapolation predicts remaining movement $(1/8)/(1-1/4)=1/6$ from $s_1$.
The iterates below $1/2$ follow the slope $3/4$ and converge to $0$, so
the actual remaining movement is $1/2$. With the valid constant $3/4$ on
the invariant interval $[0,1]$, Theorem~\ref{thm:residual} gives the
correct bounds $2$ from $s_0$ and $1/2$ from $s_1$; the latter is exact.
\end{example}

The example is one instance of a general fact.

\begin{proposition}[Histories do not determine the remaining movement]\label{prop:history}
Let $s_0>s_1>\cdots>s_N>0$ with $N\ge1$. For the feasible set $X=\{0\}$
in $B_0=[0,s_0]$, with $f\equiv0$ and $U=0$, there are two valid monotone
families with lower semicontinuous projected objectives that satisfy the
closedness condition~\eqref{eq:closed-family}, whose Jacobi iterations
from $B_0$ both produce the boxes $[0,s_k]$ for $k\le N$. In the first,
every later box equals $[0,s_N]$; in the second, the boxes converge to
$\{0\}$.
\end{proposition}
\begin{proof}
For a continuous nondecreasing function $G$ on $[0,s_0]$ with
$0\le G(s)\le s$, set $\phi_{[\ell,u]}(x)=0$ if $x\le G(u)$ and $+\infty$
otherwise, for boxes $[\ell,u]\subseteq B_0$, as in
Example~\ref{ex:not-fixed} but with continuous $G$. Validity holds
because a box containing $0$ has $\ell=0$ and $G(u)\ge0$. Monotonicity
holds because $G$ is nondecreasing, and each $\phi_B$ is the indicator of
a closed set. For~\eqref{eq:closed-family}, let boxes $C_k$ decrease to
$C$ with upper endpoints $u_k\downarrow u_\infty$, and let
$x_k\in K_0(C_k)$ converge to $x$. Then $x\in C$, and $x_k\le G(u_k)$
gives $x\le G(u_\infty)$ by continuity, so $x\in K_0(C)$. Since
$K_0([0,u])=[0,G(u)]$, the iteration is $u_{k+1}=G(u_k)$.

Let $G_1$ be the piecewise-linear interpolant of the points $(s_k,s_{k+1})$
for $0\le k<N$ together with $(0,0)$ and $(s_N,s_N)$, and let $G_2$ be the
interpolant of the same points with $(s_N,s_N)$ replaced by
$(s_N,s_N/2)$. At the increasing abscissae $0<s_N<s_{N-1}<\cdots<s_0$ the
prescribed values are nondecreasing and lie between $0$ and the abscissa,
so both functions are continuous and nondecreasing with $0\le G(s)\le s$.
Both map $s_k$ to $s_{k+1}$ for $k<N$. On $[0,s_N]$, $G_1(s)=s$ and
$G_2(s)=s/2$.
\end{proof}

After the common history, the remaining movement is $0$ in one family and
$s_N$, the largest value permitted by nesting, in the other. An upper
bound computed from $s_0,\ldots,s_N$ alone is therefore valid only if it
is at least $s_N$, the bound that nesting already gives, and no such rule
can certify a stall. This holds whether the observed ratios are close to
one or small; Section~\ref{sec:histories-con} gives explicit histories of
both kinds. The proposition concerns strictly decreasing histories. An
exact round that returns its input box, for an unchanged relaxation and
cutoff, does prove a fixed box, and nesting always bounds future movement
by the current box. The families are not McCormick relaxations. They show
that validity, monotonicity, closedness, and a finite history together do
not determine the remaining benefit; a protected box or a verified
comparison matrix must supply the missing information.

\subsection{Explicit misleading histories}\label{sec:histories-con}

Proposition~\ref{prop:history} shows that every history
$[0,s_0],\ldots,[0,s_N]$ with $s_0>\cdots>s_N>0$ is compatible both
with a stall and with convergence.
The two examples below make the two misleading patterns explicit. Both
use the construction in the proof of that proposition: for the feasible
set $\{0\}$ in $B_0=[0,1]$, with $f\equiv0$ and $U=0$, a continuous
nondecreasing map $G$ with $0\le G(s)\le s$ defines the valid monotone
family $\phi_{[\ell,s]}(x)=0$ for $x\le G(s)$ and $+\infty$ otherwise, whose
Jacobi iterates from $[0,1]$ are $[0,s_k]$ with $s_{k+1}=G(s_k)$.

\begin{example}[Ratios close to one before convergence]\label{ex:long-history-con}
Fix an integer $N\ge2$, and put $\delta=1/(2N)$ and $\sigma=1/2-\delta$.
Let
\[
 G_{\rm stall}(s)=\min\{s,\max\{\sigma,s-\delta\}\},\qquad
 G_{\rm shrink}(s)=\begin{cases}
  s/4,&0\le s\le\sigma,\\
  \sigma/4+3\sigma(s-\sigma)/(4\delta),&\sigma\le s\le1/2,\\
  s-\delta,&1/2\le s\le1.
 \end{cases}
\]
$G_{\rm stall}$ equals $s$ on $[0,\sigma]$, $\sigma$ on $[\sigma,1/2]$, and
$s-\delta$ on $[1/2,1]$; the pieces of $G_{\rm shrink}$ take the values
$\sigma/4$ at $\sigma$ and $\sigma$ at $1/2$, matching their neighbors.
Both maps are therefore continuous with nonnegative slopes, and
$0\le G(s)\le s$ holds at the endpoints of each affine piece and hence on
it. Both trajectories from $1$ subtract $\delta$ until $s_N=1/2$ and reach
$s_{N+1}=\sigma$. Every observed ratio $s_{k+1}/s_k=1-\delta/s_k$, $k\le N$,
is at least $1-1/N$. Afterwards $G_{\rm stall}$ keeps $\sigma$, while
$G_{\rm shrink}$ gives $s_{N+1+j}=4^{-j}\sigma\to0$. Arbitrarily many
ratios close to one therefore do not indicate a stall. Both maps have slope
$1$ on $[1/2,1]$, which every invariant interval containing the trajectory
includes. Any valid scalar comparison constant there is at least $1$, so
the supersolution inequality~\eqref{eq:supersolution} has no finite
solution for a positive residual, so Theorem~\ref{thm:residual} provides
no finite tail bound here.
\end{example}

\begin{example}[A small ratio before a large movement]\label{ex:small-ratio-con}
Define $G_{\rm stall}$ and $G_{\rm shrink}$ by linear interpolation of the
knots
\[
 \text{stall: }(0,0),\ (\tfrac{899}{1000},\tfrac{899}{1000}),\
  (\tfrac9{10},\tfrac{899}{1000}),\ (1,\tfrac9{10});\qquad
 \text{shrink: }(0,0),\ (\tfrac1{10},0),\ (\tfrac{899}{1000},\tfrac1{10}),\
  (\tfrac9{10},\tfrac{899}{1000}),\ (1,\tfrac9{10}).
\]
The knot values are nondecreasing and lie between $0$ and the abscissa, so
both maps are admissible. Both trajectories from $1$ begin
$1,9/10,899/1000$. The observed decrements $1/10$ and $1/1000$ have ratio
$1/100$, and geometric extrapolation predicts movement of at most
$(1/1000)/(1-1/100)=1/990$ after the first round. The stalling map stays
at $899/1000$; the shrinking map moves on to $1/10$ and then $0$, a
movement of $9/10$ after the first round.
\end{example}

Example~\ref{ex:switch-con} shows the second pattern, less extremely, in a
McCormick relaxation whose active rows change. These families do not
contradict exact certificates: nesting always bounds the remaining
movement by the current box, and an exact round that returns its input
box for an unchanged relaxation and cutoff proves a fixed box.

\subsection{Finding a comparison matrix versus checking a majorant}\label{sec:residual-cost}

Checking~\eqref{eq:supersolution} for rational $\bar d$, $M$ and $e$
takes $O(\operatorname{nnz}(M)+n)$ exact operations, and a rational
certificate $Mw\le qw$ for $\rho(M)<1$ costs the same. Neither check
establishes~\eqref{eq:matrix-lipschitz}. That hypothesis requires a bound
valid over an invariant region containing every later box, and hence over
every active regime that the support problems can enter. In
Example~\ref{ex:heron} it follows from a closed-form map. For relaxations
given by linear programs, Section~\ref{sec:constraints} derives it from
an exact cover of the region by basis regions. A derivative at one point,
a ratio of observed residuals, or a bound valid only on the current basis
region does not suffice. This paper gives no general inexpensive
construction of $M$ for rebuilt relaxations.
